\documentclass[10pt,a4paper]{amsart}
\usepackage[margin=19mm]{geometry}
\usepackage{amsmath,amssymb,mathtools}
\usepackage[scr=rsfs,cal=euler]{mathalfa}
\usepackage{xfrac}
\usepackage[unicode,hidelinks,bookmarks=false]{hyperref}
\newcommand{\ie}{i.e.\ }
\newcommand{\cf}{cf.\ }
\newcommand{\resp}{respectively\ }
\newcommand{\Subsection}{\subsection}
\providecommand{\nobreakdash}{\nobreak\hbox{-}}
\setlength{\emergencystretch}{2em}
\multlinegap0pt
\usepackage{bm}
\usepackage{bbm}
\usepackage{dsfont}
\usepackage{xspace}
\usepackage{cancel}
\usepackage{mathtools}
\usepackage{amsthm}
\makeatletter
\@ifundefined{lemma}{\newtheorem{lemma}{Lemma}}{}
\makeatother
\newcommand{\C}{\mathbb{C}}
\newcommand{\R}{\mathbb{R}}
\newcommand{\sign}{\mathrm{sign}}
\renewcommand{\S}{\mathbb{S}}
\title[A Neural Network Algorithm for KL Divergence Estimation with]{A Neural Network Algorithm for KL Divergence Estimation with Quantitative Error Bounds}
\author{Mikil Foss, Andrew Lamperski}
\date{Source: arXiv:2510.05386v1 (CC BY 4.0)}
\begin{document}
\maketitle

\noindent\emph{Source and licence.} A Neural Network Algorithm for KL Divergence Estimation with Quantitative Error Bounds. Version arXiv:2510.05386v1 (CC BY 4.0), distributed under the Creative Commons Attribution 4.0 International licence, as recorded in the arXiv licence metadata for this exact version and shown on the abstract page https://arxiv.org/abs/2510.05386v1. Full rights notice: licenses/arxiv-2510.05386-CC-BY-4.0.txt.\par\smallskip

\noindent\textit{Editorial scope.} Counted unit: the Lemma whose source label is \texttt{lem:optLinear} in the article (source0.tex lines 1140--1148, source label \texttt{lem:optLinear}), delivered together with the article's own complete original proof of it (source0.tex lines 1150--1292, one \texttt{proof} environment of 143 lines). No mathematics is written, repaired or re-derived: statement and proof are the article's own text line for line. Required context retained uncounted: lem:coordinateIntegral (lines 653--662); lem:rotationalInvariance (lines 695--701); eq:1dIntegral (lines 713--719). Every definition, display and label of those lines is retained unchanged. Omitted from this package and not redistributed: 1--652, 663--694, 702--712, 720--1139, 1149--1149, 1293--2098. Nothing inside a retained excerpt is omitted or altered: every retained body is delivered exactly as the article's lines are written, so no omission receipt of any kind accompanies this package. Citations: the retained excerpts carry no citation command at all, so no bibliography entry of the article is transcribed and no reference is redistributed. Reference adaptation: none was needed and none was made; every reference inside the delivered unit resolves inside it, so no unresolved reference or citation is produced. Printed slips kept literal: no printed word, symbol, hypothesis or number of the counted statement or of its proof was altered. Layout and class: the article's own class is \texttt{article} with packages including aistats2026, inputenc, fontenc, hyperref, url, booktabs, amsfonts, nicefrac, microtype, xcolor, float, bm, amssymb, amsmath; the wrapper is the lane's pinned \texttt{amsart} editorial wrapper at 10pt on A4, which restates the theorem-like environments the retained text needs and adds the article's own macro definitions that the retained text needs (\texttt{\textbackslash C}, \texttt{\textbackslash R}, \texttt{\textbackslash S}, \texttt{\textbackslash sign}), with extra wrapper packages (bm, bbm, dsfont, xspace, cancel, mathtools). The delivered numbering is the unit's own. Assets: no retained span contains a figure environment, a graphics-inclusion command, a PDF-page inclusion, a table or a TikZ picture; no image file is copied into this package, the package declares no asset dependency and no image file is redistributed with it. Licence: version arXiv:2510.05386v1 of this article is distributed under the Creative Commons Attribution 4.0 International licence, as recorded in the arXiv licence metadata for this exact version and shown on the abstract page https://arxiv.org/abs/2510.05386v1. A full rights notice is delivered at licenses/arxiv-2510.05386-CC-BY-4.0.txt. Characters: every retained excerpt is pure ASCII, so the delivered engine, XeLaTeX with its default Latin Modern text font, reports no missing character.
\begin{lemma}
\label{lem:coordinateIntegral}
If $f\in L^1(\S^{n-1})$, then its integral can be expressed in the following equivalent ways:
\begin{align*}
\int_{\S^{n-1}}f(w)\mu_{n-1}(dw)&=\int_{\Phi}
f(h(\phi))\sqrt{\det(Dh(\phi)^\top Dh(\phi))}d\phi \\
&=\int_{\Phi}f(h(\phi))\left(\prod_{i=1}^{n-2}\sin^{n-1
-i}(\phi_i)\right)d\phi.
\end{align*}
\end{lemma}

\begin{lemma}
\label{lem:rotationalInvariance}
If $f:\S^{n-1}\to \C$ is in $L^1(\S^{n-1})$ and $U$ is an $n\times n$ orthogonal matrix, then
$$
\int_{\S^{n-1}}f(w)\mu_{n-1}(dw)=\int_{\S^{n-1}}f(Uz)\mu_{n-1}(dz).
$$
\end{lemma}

\begin{lemma}
\label{eq:1dIntegral}
If $n\ge 3$,  $g\in L^1(\R)$ and $v\in \R^n$, then
$$
\int_{\S^{n-1}}g(v^\top w)\mu_{n-1}(dw)=A_{n-2}\int_0^\pi g(\|v\|_2 \cos(\phi_1))\sin(\phi_1)^{n-2}d\phi_1.
$$
\end{lemma}

\begin{lemma}
\label{lem:optLinear}
For $n\ge 1$ and $v\in \R^n$,
the function $q:\S^{n-1}\to \R$ defined by $q(w)=\frac{\|v\|_2}{H_{n-1}}\sign(v^\top w)$, where $H_{n-1}=\frac{2\pi^{\frac{n-1}{2}}}{\Gamma\left(\frac{n+1}{2}\right)}$, is an optimal solution to the following functional optimization problem:
\begin{align*}
&\min_f && \|f\|_{L^{\infty}(\S^{n-1})}\\
&\textrm{subject to} && \int_{\S^{n-1}}w f(w)\mu_{n-1}(dw)=v.
\end{align*}
\end{lemma}

\begin{proof}
For $v=0$, the function becomes $q(w)=0$ for all $w\in \S^{n-1}$, which is feasible and achieves the smallest possible norm. Thus the lemma holds in this case. The rest of the proof will focus on the case that $v\ne 0$.

The proof proceeds as follows. We show that $q(w)=\frac{\|v\|_2}{H_{n-1}}\sign(v^\top w)$ is feasible. Note  here that the value obtained is $\|v\|_2/H_{n-1}.$ Then we will construct the Lagrange dual to the optimization problem and find a dual solution also obtaining value $\|v\|_2/H_{n-1}.$ It then will follow from weak duality that $q$ is optimal. 

To show that $q$ is feasible, we must show that the constraint holds. For this calculation, it is more convenient to work in coordinates in which $v$ is aligned with the first unit vector. To this end, set $u_1=v/\|v\|_2$, and let $U=\begin{bmatrix}u_1 & \ldots & u_n\end{bmatrix}$ be an orthogonal matrix. Let $z=U^\top w$. Then, using Lemma~\ref{lem:rotationalInvariance} on rotational invariance of the sphere, $q$ satisfies the constraint if and only if:
\begin{align*}
\int_{\S^{n-1}}U^\top w q(w)\mu_{n-1}(dw)&=
\frac{\|v\|_2}{H_{n-1}}\int_{\S^{n-1}}z\sign(z_1)\mu_{n-1}(dz)\\
&=\frac{\|v\|_2}{H_{n-1}}\int_{\S^{n-1}}
\begin{bmatrix}
|z_1|\\
z_2 \sign(z_1) \\
\vdots \\
z_n\sign(z_1)
\end{bmatrix}
\mu_{n-1}(dz) \\
&=U^\top v = \begin{bmatrix}
\|v\|_2 \\
0 \\
\vdots \\
0
\end{bmatrix}.
\end{align*}

Thus, it suffices to show that 
\begin{align}
\label{eq:coordinateIntegral}
H_{n-1}&=\int_{\S^{n-1}}|z_1|\mu_{n-1}(dz) = \frac{2\pi^{n-1}}{\Gamma\left(\frac{n+1
}{2}\right)} \\
\label{eq:signOrthogonal}
0&=\int_{\S^{n-1}}z_i\sign(z_1)\mu_{n-1}(dz) \textrm{ for } i=2,\ldots,n.
\end{align}

For $n=1$, only \eqref{eq:coordinateIntegral}  must be shown, and in this case, both sides evaluate to $2$. For $n=2$, both sides of \eqref{eq:coordinateIntegral} evaluate to $4$, and \eqref{eq:signOrthogonal} holds by direct calculation. 

For $n\ge 3$, Lemma~\ref{eq:1dIntegral}, followed by  some manipulations gives:
\begin{align*}
\int_{\S^{n-1}}|z_1|\mu_{n-1}(dz)&=A_{n-2}\int_0^\pi |\cos(\phi_1)|\sin^{n-2}(\phi_1)d\phi \\
&=2A_{n-2}\int_0^{\pi/2} \cos(\phi_1)\sin^{n-2}(\phi_1)d\phi \\
&=\frac{2A_{n-2}}{n-1} \\
&=\frac{4\pi^{\frac{n-1}{2}}}{(n-1)\Gamma\left(\frac{n-1}{2}\right)}\\
&=\frac{2\pi^{\frac{n-1}{2}}}{\Gamma\left(\frac{n+1}{2}\right)}.
\end{align*}
Thus, \eqref{eq:coordinateIntegral} holds for all $n\ge 1$.

Now we evaluate the integrals from \eqref{eq:signOrthogonal} via Lemma~\ref{lem:coordinateIntegral}. %For $i=n-1,n$, integrating over $\phi_1,\ldots,\phi_{n-2}$ gives:
%\begin{align*}
%\int_{\S^{n-1}}z_{n-1}\sign(z_1)\mu_{n-1}(dz)&\propto \int_0^{2\pi}\sin(\phi_{n-1})d\phi_{n-1}=0 \\
%\int_{\S^{n-1}}z_{n}\sign(z_1)\mu_{n-1}(dz)&\propto \int_0^{2\pi}\cos(\phi_{n-1})d\phi_{n-1}=0.
%\end{align*}
%
For $i=2,\ldots,n$, integrating over $\phi_2, \ldots, \phi_{n-1}$ gives:
\begin{align*}
\int_{\S^{n-1}}z_i\sign(z_1)\mu_{n-1}(dz)
&\propto \int_{0}^\pi \sign(\cos(\phi_1))\sin^{n-1}(\phi_1)d\phi_1 \\
&=\int_0^{\pi/2}\sin^{n-1}(\phi_1)d\phi_1 - \int_{\pi/2}^{\pi}\sin^{n-1}(\phi_1)d\phi_1 = 0.
\end{align*}
The final equality follows from substitution $\psi=\pi-\phi_1$ in the second integral:
\begin{align*}
\int_{\pi/2}^{\pi}\sin^{n-1}(\phi_1)d\phi_1 &=\int_0^{\pi/2}\sin^{n-1}(\pi-\psi)d\psi \\
&=\int_0^{\pi/2}\sin^{n-1}(\psi)d\psi. 
\end{align*}
Thus, \eqref{eq:signOrthogonal} holds for all $n\ge 1$ and all $2\le i \le n$.

Since \eqref{eq:coordinateIntegral} and \eqref{eq:signOrthogonal} hold, the function $q$ is feasible, giving objective value $\|v\|_2/H_{n-1}.$

Now we will derive the Lagrange dual, and find a dual solution obtaining value $\|v\|_2/H_{n-1}$.

The optimization problem from the lemma statement can be posed equivalently as:
\begin{align*}
&\min_{t,f} && t\\
&\textrm{subject to} && \int_{\S^{n-1}}w f(w)\mu_{n-1}(dw)=v \\
&&& -t\le f(w)\le t \textrm{ for almost all } w\in\S^{n-1}.
\end{align*}
This is an infinite dimensional linear program over variables $(t,f)\in \R\times L^{\infty}(\S^{n-1})$.

The Lagrangian of the reformulated problem is given by
\begin{align*}
%\MoveEqLeft
L(t,f,\lambda,\alpha,\beta)%\\
&=t+\lambda^\top \left(v- \int_{\S^{n-1}}w f(w)\mu_{n-1}(dw)\right)
\\ &+\int_{\S^{n-1}}\left(-\alpha(w)(t+f(w))+\beta(w)(f(w)-t) \right)\mu_{n-1}(dw) \\
&=\lambda^\top v+
t\left(1-\int_{\S^{n-1}}(\alpha(w)+\beta(w))\mu_{n-1}(dw) \right)\\
&+\int_{\S^{n-1}}f(w)\left(\beta(w)-\alpha(w)-\lambda^\top w \right)\mu_{n-1}(dw),
\end{align*}
with dual variables $(\lambda,\alpha,\beta)\in \R^n\times L^1(\S^{n-1})\times L^1(\S^{n-1}).$

The associated dual problem is given by:
\begin{align*}
&\max_{\lambda,\alpha,\beta} && \lambda^\top v \\
&\textrm{subject to} && \alpha(w)\ge 0, \beta(w)\ge 0, \textrm{ for almost all } w\in \S^{n-1} \\
&&& \int_{\S^{n-1}}(\alpha(w)+\beta(w))\mu_{n-1}(dw)=1\\
&&& \beta(w)-\alpha(w)=\lambda^\top w, \textrm{ for almost all } w\in \S^{n-1}.
\end{align*}

We claim that the dual problem is equivalent to:
\begin{subequations}
\label{eq:vDual}
\begin{align}
&\max_{\lambda,\alpha,\beta} && \lambda^\top v \\
&\textrm{subject to} && \int_{\S^{n-1}}|\lambda^\top w|\mu_{n-1}(dw)\le 1.
\end{align}
\end{subequations}

Indeed, for a dual feasible $(\lambda,\alpha,\beta)$, we must have $|\lambda^\top w|\le \alpha(w)+\beta(w)$, and so $\lambda$ is feasible for \eqref{eq:vDual}.

Conversely, given any $\lambda$ feasible for \eqref{eq:vDual}, let 
$$
s=\int_{\S^{n-1}}|\lambda^\top w|\mu_{n-1}(dw)\le 1.
$$
Then, we can construct corresponding dual feasible $\alpha$ and $\beta$ by setting
\begin{align*}
\alpha(w)&=\begin{cases}
-\lambda^\top w +\frac{1-s}{2A_{n-1}}& \lambda^\top w <0 \\
\frac{1-s}{A_{n-1}} & \lambda^\top w \ge 0
\end{cases}\\
\beta(w)&=\begin{cases}
\lambda^\top w + \frac{1-s}{2A_{n-1}} &\lambda^\top w > 0 \\
\frac{1-s}{2A_{n-1}}& \lambda^\top w \le 0.
\end{cases}
\end{align*}

Recall that we are examining the case that $v\ne 0.$
Let 
$$
\lambda = \frac{1}{\|v\|_2 H_{n-1}}v.
$$
Note that $\lambda^\top v = \|v\|_2/H_{n-1}$, which was the value obtained by $q$ on the primal problem. As discussed above, $\lambda$ will correspond to a dual feasible solution as long as it is feasible for \eqref{eq:vDual}.

Recall the change of coordinates from above, $z=U^\top w$, where $z_1=v^\top w/\|v\|_2$. Using rotational invariance of the sphere, Lemma~\ref{lem:rotationalInvariance},
gives:
\begin{align*}
\int_{\S^{n-1}}|\lambda^\top w|\mu_{n-1}(dw)&=\frac{1}{\|v\|_2 H_{n-1}}\int_{\S^{n-1}}|v^\top w|\mu_{n-1}(dw)\\
&=\frac{1}{H_{n-1}}\int_{\S^{n-1}}|z_1|\mu_{n-1}(dz) \\
&=1,
\end{align*}
where the final equality is from \eqref{eq:coordinateIntegral}.

Thus, $\lambda$ is feasible for \eqref{eq:vDual}. By weak duality, the value achieved, $\lambda^\top v=\|v\|_2/H_{n-1}$,
is a lower bound on the achievable value for the optimization problem from the lemma statement. Thus, $q$ must be  optimal. 
\end{proof}

\end{document}
